% Введение. Структура по ВКР для 01.03.02.
% Источники фактов: knowledge-base/.

\thesisChapterStar{Введение}

\textbf{Актуальность темы.} Цифровизация образовательного процесса в российских вузах остаётся неравномерной: при наличии массовых решений для отчётности и контингента, повседневный инструмент --- учебное расписание --- во многих учебных заведениях ведётся вручную в Excel-таблицах и публикуется в виде PDF-файлов на сайтах вузов. Это создаёт системные потери на стороне администрации и неудобство на стороне студентов и преподавателей. Запрос рынка на современное, мобильное и интегрированное решение подтверждается как наличием прямых конкурентов из числа тяжёлых ERP-систем, так и распространением неофициальных приложений-парсеров.

\textbf{Цель работы.} Спроектировать и реализовать комплекс клиентских приложений (iOS, Android, Web) SaaS-платформы <<ВКурсе>>\footnote{В магазинах приложений (App~Store, Google~Play, RuStore) клиентские приложения опубликованы под полным маркетинговым названием <<ВКурсе: Расписание занятий>> (англ.~\emph{VCourse: Class schedule}), используемым для повышения находимости по поисковым запросам.}, обеспечивающий единый пользовательский опыт работы с учебным расписанием для студентов, преподавателей и сотрудников учебных заведений.

\textbf{Задачи работы.}
\begin{enumerate}
  \item Проанализировать существующие решения для управления учебным расписанием (прямых и непрямых конкурентов), выявить ограничения и сформировать требования к клиентской части платформы.
  \item Спроектировать единый продуктовый подход с учётом нативных рекомендаций каждой платформы (Apple Human Interface Guidelines, Material Design 3 Expressive, shadcn/ui).
  \item Разработать клиентские мобильные приложения для iOS и Android с поддержкой офлайн-режима и персонализации.
  \item Разработать веб-приложение, включающее публичную часть (лендинг и просмотр расписания) и закрытый сервис администрирования учебного заведения.
  \item Провести апробацию результатов: публикация в магазинах приложений, сбор метрик использования, организация диалога с учебными заведениями.
\end{enumerate}

\textbf{Объект исследования} --- процесс управления и доставки учебного расписания в образовательной организации.

\textbf{Предмет исследования} --- клиентская часть SaaS-платформы для управления учебным расписанием.

\textbf{Методология.} Работа сочетает методы анализа предметной области (сравнительный обзор существующих систем), UX/UI-проектирования (раздельная дизайн-система под каждую целевую платформу), итеративной разработки клиентских приложений и опытной апробации (публикация в магазинах приложений, сбор пользовательской обратной связи и количественных метрик использования).

\textbf{Научная новизна и практическая значимость.} В работе предложен и реализован комплексный подход к построению клиентской части SaaS-платформы для управления учебным расписанием, отличающийся от существующих российских решений сочетанием SaaS-модели поставки, нативной мобильной составляющей и независимости от сторонних ERP-экосистем. Практическая значимость подтверждается публикацией клиентских приложений в App Store, Google Play и RuStore, ежедневной активной аудиторией порядка 1000~уникальных пользователей и начатым диалогом с учебными заведениями о пилотном внедрении.

\textbf{Апробация результатов.} Клиентские приложения опубликованы в App Store (сентябрь 2025~г.), Google Play и RuStore (февраль 2026~г.). Ход разработки и обновления освещаются в Telegram-канале проекта (более 410 подписчиков). Проект представлен инвестиционному фонду <<ТРАМПЛИН>> (Омск, апрель 2026~г.) и Фонду содействия инновациям в рамках гранта <<Студенческий стартап>>. В мае 2026 года состоялась первая очная встреча с потенциальным пилотным партнёром --- Омским государственным аграрным университетом.

\textbf{Личный вклад автора.} Автором выполнены анализ предметной области, продуктовое и UX/UI-проектирование всех клиентских компонентов экосистемы, полная реализация мобильных приложений для iOS и Android, веб-лендинга, веб-приложения просмотра расписания и веб-сервиса администрирования, подготовка и проведение релизов в магазинах приложений, а также организационная работа с учебными заведениями. Серверная часть платформы разработана соавтором в рамках отдельной выпускной квалификационной работы; в настоящей работе она рассматривается только в объёме интеграционного контракта с клиентской частью.

\textbf{Структура работы.} Работа выполнена по направлению подготовки 01.03.02 <<Прикладная математика и информатика>> в соответствии с требованиями федерального государственного образовательного стандарта высшего образования~\cite{fgos-01-03-02}; библиографические ссылки оформлены по ГОСТ Р 7.0.5-2008~\cite{gost-7-0-5-2008}. Работа состоит из введения, трёх глав, заключения, списка литературы и приложений. В первой главе проводится анализ предметной области и формулируются требования к клиентской части. Во второй главе описывается проектирование архитектуры экосистемы, пользовательского опыта и дизайн-систем под каждую платформу. В третьей главе раскрывается реализация iOS-, Android- и веб-приложений, а также апробация результатов.
